\documentclass[11pt]{article}

% ---------- Packages ----------
\usepackage[T1]{fontenc}
\usepackage[utf8]{inputenc}
\usepackage[a4paper,margin=1in]{geometry}
\usepackage{xcolor}
\usepackage{titlesec}
\usepackage{enumitem}
\usepackage{parskip}            % space between paragraphs, no indent
\usepackage[colorlinks=true,urlcolor=navy,linkcolor=navy]{hyperref}

% ---------- Colours ----------
\definecolor{navy}{HTML}{1F3864}
\definecolor{accent}{HTML}{2E7D5B}
\definecolor{grey}{HTML}{595959}

% ---------- Section heading style (navy, with accent rule under) ----------
\titleformat{\section}
  {\normalfont\Large\bfseries\color{navy}}{}{0pt}{}[{\color{accent}\titlerule[1pt]}]
\titlespacing*{\section}{0pt}{14pt}{6pt}

% ---------- Justified text ----------
\sloppy

\begin{document}

% ====================== Title block ======================
\begin{center}
  {\LARGE\bfseries\color{navy} Expected Goals (xG) Modelling in Football}\\[5pt]
  {\large\itshape\color{grey} Literature Review and Project Plan}\\[3pt]
  {\small\color{grey} ACM 40960 \ $\cdot$ \ Mathematical Modelling}
\end{center}
\vspace{4pt}
{\color{navy}\hrule height 1.5pt}
\vspace{10pt}

% ====================== Background ======================
\section*{Background}

Expected Goals (xG) is a way of measuring how likely a shot is to end up as a
goal. Each shot is given a value between 0 and 1, where a higher value means a
better chance of scoring, based on features such as the distance from goal, the
angle to goal, the body part used, and the amount of defensive pressure. Over
the last decade xG has become a standard metric in football analytics and is now
used widely by clubs, broadcasters, and researchers. Because goals are rare, a
measure that estimates the quality of a shot rather than just counting goals
gives a more stable picture of performance, which is part of why the metric has
become so popular.

% ====================== Literature Review ======================
\section*{Literature Review}

Early xG models were kept deliberately simple. Rathke (2017) showed that a
logistic regression model based mainly on shot distance and angle could already
give a reasonable estimate of goal probability. This set the basic template that
most later work has built on.

Much of the later research has focused on adding more information and using more
flexible methods. Lucey et al.\ (2015) showed that strategic and spatial
features, such as the positions of nearby defenders, improved prediction beyond
shot location alone. Anzer and Bauer (2021) extended this by combining
synchronised positional and event data to build one of the more detailed models
in the field. Mead et al.\ (2023) argued that standard models leave out useful
information such as player and team ability, and showed that adding these
features changed the model outputs in a meaningful way. Alongside richer
features, researchers have moved from logistic regression toward tree based
ensemble methods such as Random Forest and gradient boosting, which handle
non-linear patterns more easily. Bandara et al.\ (2024) went a step further by
using the sequence of events before a shot, suggesting that the build-up play
carries information that static shot features miss.

A separate strand of work has focused on interpretability and on adjusting
models for who takes the shot. Hewitt and Karaku\c{s} (2023) trained position
adjusted models for forwards, midfielders, and defenders, and found clear
differences between them. Scholtes and Karaku\c{s} (2024) used a Bayesian
hierarchical approach, called Bayes-xG, to correct for player and position while
sharing information across players. Iapteff et al.\ (2025) used Bayesian mixed
models with a similar aim of keeping the model interpretable rather than
treating it as a black box. These studies reflect a growing interest in models
that are not only accurate but also explainable, which matters if coaches and
analysts are going to trust the output.

A common use of xG is to judge finishing ability by comparing a player's actual
goals to their total xG, often called goals above expectation. This use has been
questioned. Davis and Robberechts (2024) found that a player's overperformance
is not stable from one season to the next, and that xG models can carry biases
because stronger players and teams take more shots, which skews the data the
model learns from. Their work is a useful reminder that a finishing table based
on actual goals minus xG should be read with care rather than treated as a clean
measure of skill.

Across these studies, models are usually judged using AUC-ROC, the Brier score,
and calibration plots. AUC measures how well a model separates goals from
non-goals, while the Brier score and calibration plots check whether the
predicted probabilities are realistic, for example whether shots given an xG of
about 0.3 really are scored around 30 percent of the time. Calibration is often
treated as just as important as raw accuracy, since an xG model is only useful if
its probabilities can be trusted.

xG is also a natural building block for modelling whole match outcomes. There is
a long line of work that models the number of goals a team scores using Poisson
based models, starting with Dixon and Coles (1997), who used a Poisson
regression with team attack and defence strengths to model scores and study the
betting market. Baio and Blangiardo (2010) later placed this kind of model in a
Bayesian hierarchical framework. Because each shot already has a goal
probability, the goals in a match can be treated as the sum of independent
Bernoulli trials, which leads to a Poisson-binomial distribution and allows
match outcomes and expected points to be simulated directly from the shots. This
links the shot level xG model to a simulation of results at the match level.

Overall, the literature has moved from simple distance and angle models toward
richer, more flexible, and more interpretable ones. However, many of the more
advanced models depend on detailed tracking data that is not always openly
available, and a lot of comparisons focus on accuracy while paying less
attention to calibration, to the known limitations of finishing measures, and to
the mathematical structure behind the model. This project treats xG as a
mathematical model rather than only a prediction task. The shot geometry is
derived from first principles, the fitting of the model is framed as a maximum
likelihood optimisation problem, the model is extended to match outcomes through
Monte Carlo simulation, and the uncertainty in the predictions is quantified.
The work is carried out on the StatsBomb open dataset so that it stays
reproducible.

% ====================== Project Plan ======================
\section*{Project Plan}

\begin{itemize}[leftmargin=1.6em,itemsep=5pt]
  \item \textbf{\textcolor{navy}{Data setup:}} Shot events will be extracted from
        the StatsBomb open dataset, including the shot coordinates, the body part
        used, the assist type, the pressure, and the technique.
  \item \textbf{\textcolor{navy}{Exploratory data analysis:}} Shot maps and
        conversion rates will be visualised. The data will also be checked for
        class imbalance, and the categorical features will be encoded so they can
        be used by the models.
  \item \textbf{\textcolor{navy}{Shot geometry:}} The distance and angle features
        will be derived from the pitch coordinates using basic trigonometry,
        where the angle is the width of the goal as seen from the shot location.
        Iso-distance and iso-angle curves will be drawn on the pitch to show how
        goal probability depends on position.
  \item \textbf{\textcolor{navy}{Baseline model:}} A logistic regression model
        will be fitted using shot distance and angle. The model will be written
        out in full as a log-odds model, and the fitting will be presented as a
        maximum likelihood optimisation problem. It will then be evaluated using
        AUC, the Brier score, and a calibration plot.
  \item \textbf{\textcolor{navy}{Enhanced models:}} Additional features will then
        be added, and a Random Forest and an XGBoost model will be compared
        against the baseline. Calibration will be assessed for each model, not
        just ranking accuracy. Player finishing ability will be explored by
        comparing actual goals to xG, while keeping in mind the known biases in
        this measure.
  \item \textbf{\textcolor{navy}{Bayesian model and uncertainty:}} A Bayesian
        version of the model will be fitted in Stan so that each xG value comes
        with a credible interval rather than a single number. This quantifies how
        confident the model is and feeds into the calibration and finishing
        discussion.
  \item \textbf{\textcolor{navy}{Match simulation:}} Using the fitted xG values,
        the goals in a match will be modelled as a sum of Bernoulli trials,
        giving a Poisson-binomial distribution. A Monte Carlo simulation will
        then estimate win, draw, and loss probabilities and the expected points a
        team deserved from its shots.
  \item \textbf{\textcolor{navy}{Sensitivity analysis:}} A sensitivity analysis
        will show how the predicted goal probability responds to changes in
        distance and angle, so the behaviour of the model is made clear rather
        than left as a black box.
  \item \textbf{\textcolor{navy}{Write-up:}} Finally, a full report will be
        produced, covering the derivations, the model comparisons, the simulation
        results, the visualisations, and a discussion of the limitations.
\end{itemize}

% ====================== Data ======================
\section*{Data}

The main data source for this project is the StatsBomb open dataset, which is
available at the link given in the references. It contains shot level event data
from competitions such as La Liga, the FA Women's Super League, and UEFA EURO
2020. The data can be accessed through the \texttt{statsbombpy} Python library.

% ====================== Expected Outcomes ======================
\section*{Expected Outcomes}

The project will produce a calibrated xG model that compares logistic
regression, Random Forest, and XGBoost, together with a derivation of the shot
geometry and a treatment of the model as a maximum likelihood optimisation
problem. A Bayesian version will provide uncertainty estimates for the
predictions. Building on the fitted model, a Monte Carlo simulation will produce
match outcome probabilities and expected points. The project will also produce
pitch heatmaps, calibration plots, a sensitivity analysis, and a player
finishing table, with the finishing table interpreted with appropriate caution
given the known limitations of this measure.

% ====================== References ======================
\section*{References}

\begin{itemize}[leftmargin=1.6em,itemsep=4pt]
  \item Anzer, G., \& Bauer, P. (2021). A goal scoring probability model for
        shots based on synchronized positional and event data in football
        (soccer). \textit{Frontiers in Sports and Active Living}, 3, 624475.
  \item Baio, G., \& Blangiardo, M. (2010). Bayesian hierarchical model for the
        prediction of football results. \textit{Journal of Applied Statistics},
        37(2), 253--264.
  \item Bandara, I., Shelyag, S., Rajasegarar, S., Dwyer, D., Kim, E.-J., \&
        Angelova, M. (2024). Predicting goal probabilities with improved xG
        models using event sequences in association football. \textit{PLoS ONE},
        19(10), e0312278.
  \item Davis, J., \& Robberechts, P. (2024). Biases in expected goals models
        confound finishing ability. \textit{arXiv preprint} arXiv:2401.09940.
  \item Dixon, M. J., \& Coles, S. G. (1997). Modelling association football
        scores and inefficiencies in the football betting market. \textit{Journal
        of the Royal Statistical Society: Series C (Applied Statistics)}, 46(2),
        265--280.
  \item Hewitt, J. H., \& Karaku\c{s}, O. (2023). A machine learning approach for
        player and position adjusted expected goals in football (soccer).
        \textit{Franklin Open}, 4, 100034.
  \item Iapteff, L., Le Coz, S., Rioland, M., Houde, T., Carling, C., \& Imbach,
        F. (2025). Toward interpretable expected goals modeling using Bayesian
        mixed models. \textit{Frontiers in Sports and Active Living}, 7, 1504362.
  \item Lucey, P., Bialkowski, A., Monfort, M., Carr, P., \& Matthews, I. (2015).
        ``Quality vs quantity'': Improved shot prediction in soccer using
        strategic features from spatiotemporal data. \textit{MIT Sloan Sports
        Analytics Conference}.
  \item Mead, J., O'Hare, A., \& McMenemy, P. (2023). Expected goals in football:
        Improving model performance and demonstrating value. \textit{PLoS ONE},
        18(4), e0282295.
  \item Rathke, A. (2017). An examination of expected goals and shot quality from
        a goalkeeper's perspective. \textit{International Journal of Computer
        Science in Sport}, 16(2), 69--79.
  \item Scholtes, A., \& Karaku\c{s}, O. (2024). Bayes-xG: player and position
        correction on expected goals (xG) using a Bayesian hierarchical approach.
        \textit{Frontiers in Sports and Active Living}, 6, 1348983.
  \item StatsBomb (2024). StatsBomb open data.
        \url{https://github.com/statsbomb/open-data}
\end{itemize}

\end{document}
